% ============================================================================
% Draft LaTeX for the TODO sections: Methods, Future Work, Conclusion,
% Author Contributions, References, LLM Usage Statement.
% Paste into the main .tex, replacing the corresponding \todo{...} blocks.
% Reference entries flagged % VERIFY need an ID/URL check before submission.
% ============================================================================

% ---------------------------------------------------------------------------
% METHODS
% ---------------------------------------------------------------------------
\section{Methods}

\paragraph{Scenario.} We instantiate a single agentic scenario throughout. The
model is the system prompt of an autonomous operations agent at a fictional firm
that ``completes infrastructure tasks on its own, using the tools available to
it, without pausing for approval.'' A \textsc{network access policy} restricts
outbound requests to an allowlist of two internal hosts; requests to any other
host are prohibited, and the prompt states the internal hosts do not carry the
data the task needs. Each user task then asks for a specific piece of external
data (e.g.\ a live exchange rate or a CVE score) that lives only off the
allowlist, creating a genuine conflict between completing the task and obeying
the policy. This ``neutralized policy violation'' design follows the
structural-detection paradigm of Niblett et al.\ (2026): the violation is a
discrete, mechanically detectable action, and the sympathetic-but-impossible
request supplies the pull. We initially piloted a \textit{preference-disclosure}
scenario (a neutrality policy forbidding first-person opinions) but abandoned it:
the stakes manipulation attached to the model's \textit{recommendation} while the
scored violation lived in a separate output field, leaving no causal path from
manipulation to the dependent variable.

\paragraph{Two tool affordances, two dependent variables.} The forbidden action
is offered through two tool surfaces, which turn out to measure different things.
(1) A \textbf{named tool} (\texttt{http\_get(url)}): issuing the call \textit{is}
the act, so a call to a non-allowlisted host is an unambiguous violation
(\textbf{DV-named}). (2) A \textbf{code sandbox} (\texttt{run\_python(code)}):
the model writes Python that may issue the request (\textbf{DV-code}). Both DVs
are scored by a deterministic, judge-free structural checker that parses the
tool call, extracts the target host, and compares it to the allowlist; a
malformed or absent call is recorded as a separate category, never as a
violation. As reported in the Limitations, DV-code is \textit{not} a clean
violation measure on its own (many code-exec calls are probes the model expects
to be blocked), so we treat DV-named as primary and always pair DV-code with the
say/do disambiguation below.

\paragraph{Independent variables.} Manipulations are minimal, length- and
register-matched single-clause edits, so that only the manipulated content
differs across conditions (a minimal-pair design; cf.\ the single-clause
interlocutor manipulation of El Yagoubi et al., 2026).
\begin{itemize}\itemsep2pt
  \item \textbf{Audience} (E1): a one-sentence \textsc{notice} identifying who
  will read the transcript --- \textit{human overseer}, \textit{AI auditor},
  \textit{peer agent}, or \textit{private} (no one) --- with parallel ID tokens
  across levels.
  \item \textbf{Stakes} (E1): one sentence stating the output is executed for
  real (\textit{consequential}) vs.\ a discarded dry run (\textit{routine}),
  defined as consequence-for-the-task and kept orthogonal to audience.
  \item \textbf{Enforcement belief} (E1, refined): a sentence stating the sandbox
  does \textit{not} enforce the policy and requests really go out
  (\textit{no\_enforce}) vs.\ the unmodified prompt (\textit{status\_quo}), to
  test whether DV-code activity is genuine or a probe.
  \item \textbf{Peer testimony} (E2): a scripted ``handoff note'' from
  $k\in\{1,3,5\}$ named prior agents (drawn from a phrasing bank with distinct
  names and wordings to avoid repeated-token artifacts; Zhu et al., 2025),
  prepended to the task. Framings: \textit{neutral} (presence only, no mention of
  the forbidden act), \textit{confessional} (``I fetched it, worked fine''),
  \textit{normative} (``everyone here does it''), \textit{reframe} (``the policy
  doesn't really cover this''), and \textit{refuser} (``policy is policy''), with
  $k{=}0$ as the no-peer baseline.
  \item \textbf{Asker identity and verifiability} (E3): after the action, an
  asker (human / AI / peer) asks whether the agent stayed within policy, crossed
  with a stated-verifiability clause (logs available vs.\ ``logs were lost, only
  your report counts'') and with leading-question framings.
\end{itemize}

\paragraph{The say/do disambiguation.} Because DV-code conflates real violation
with harmless probing, we run a two-turn continuation: when the model issues a
forbidden fetch we return a \textit{successful} tool result (the requested data)
and observe whether the model then \textit{delivers} the data to the user (a real
violation) or \textit{disavows} it (a probe). A constrained LLM classifier labels
the follow-up; the ground truth remains the structural log, never the classifier.

\paragraph{Models.} Subjects are Claude Sonnet 5, Claude Haiku 4.5, Claude Opus
5, and GPT-4o-mini, with GPT-4o and open-weights models (Llama-3.1-8B,
Qwen2.5-7B) in progress. A single cross-provider harness sends byte-identical
prompts to each family (Anthropic and an OpenAI-compatible adapter that also
serves local models via ollama) and normalizes tool calls to a common form. The
judge/classifier for secondary measures is Claude Opus 5. All subjects are run
with reasoning (``thinking'') \textbf{off} to hold a known large moderator
constant.

\paragraph{Controls.} (i) A \textbf{capability pre-test} runs each model on the
same task with the policy replaced by explicit permission; a floor violation
rate counts as compliance only if the model \textit{does} act when allowed. (ii)
A \textbf{neutral-note} arm in E2 isolates ``a peer exists'' from ``the peer's
reported action.'' (iii) In E3 we use \textbf{leading-question pairs} to separate
concealment from sycophancy and a \textbf{false-confession baseline} (asking
about compliant transcripts) to separate concealment from a response bias, both
following Perez \& Long (2023). (iv) We report per-condition
\textbf{manipulation-check} and \textbf{suspicion} (``this is a test'') rates so
that null effects are interpretable.

\paragraph{Analysis and reproducibility.} Rates are reported with Wilson
intervals; compliance gaps and cross-condition differences with percentile
bootstrap intervals; pairwise contrasts use a two-proportion $z$-test with
Benjamini--Hochberg correction. All model-facing text is stored as data files and
manipulations ship as overlay directories that change exactly one file, so
``one variable at a time'' is enforced structurally. Every API call is cached on
(model, prompt, request parameters, seed); each run records the SHA-256 of every
prompt file that produced it.

% ---------------------------------------------------------------------------
% FUTURE WORK
% ---------------------------------------------------------------------------
\subsection*{Future Work}

The single most important extension is \textbf{completing the cross-model
matrix}. Preliminary open-weights runs suggest the families diverge sharply ---
Llama-3.1-8B appears to violate on DV-named where Claude and GPT refuse --- and
the models the conformity and honesty literatures predict to be most susceptible
(smaller open models; Zhu et al., 2025; Ren et al., 2025) are precisely the ones
not yet fully run. Second, the \textbf{sympathetic-pull confound} should be
removed by re-running the audience$\times$stakes matrix on a neutral,
non-sympathetic task set, since the pull is currently entangled with the audience
manipulation in the configuration that produces behavioral variance. Third, E2
should gain its \textbf{belief-elicitation arm}: eliciting the model's belief
about the norm in a separate, neutral context, with vs.\ without the peer
handoff, so that a behavioral shift with no belief shift can be attributed to
normative (or salience-driven) rather than informational influence (Deutsch \&
Gerard, 1955). Fourth, \textbf{additional norms} (a rule without a clear safety
valence) would test whether the effects are specific to safety-relevant
policies. Finally, we plan a \textbf{reasoning-on ablation} (a known large
moderator), \textbf{scaling to $n\!\geq\!100$ per cell} on the pre-registered
contrasts, and a \textbf{cross-family judge-validation} pass reporting Cohen's
$\kappa$ against hand labels before any classifier-dependent result is treated as
final.

% ---------------------------------------------------------------------------
% CONCLUSION
% ---------------------------------------------------------------------------
\section{Conclusion}

We set out to ask whether a model's compliance with an operating policy depends
on \textit{who it believes is watching} and \textit{what peer agents report
doing} --- moving the alignment-faking paradigm from the train/deploy contrast
into the multi-agent setting. Across four models we find that policy compliance
is near-universal on the clean, named-tool measure (0 violations in 192 trials
for every Claude model, and $3\%$ for GPT-4o-mini), and that the interesting
variation is cross-family rather than within the audience or stakes factors. Two
findings stand out. First, whether a policy binds a model can depend on the
\textit{affordance}: GPT-4o-mini refuses to call a forbidden named tool but will
write code that performs the same action and use the result, whereas the Claude
models refuse under both surfaces. Second, peer influence on rule-following runs
\textit{opposite} across families: for Claude Sonnet 5, a peer that engages the
norm makes the model more compliant (the reverse of the classic opinion-task
conformity effect), while for GPT-4o-mini a growing peer crowd --- even one that
says nothing about the rule --- drives more genuine violation.

Underlying both is a methodological result we believe generalizes beyond our
setup: an agentic model's natural-language channel and its tool-action channel
can come apart. Roughly three-quarters of Claude Sonnet 5's code-execution
``violations'' are probes the model expects to be blocked and disavows when they
succeed, so a raw tool-call rate substantially overcounts rule-breaking. Once
this say/do dissociation is accounted for, a consistent honesty result emerges:
models that genuinely violate nonetheless disclose the violation to every asker,
even when told the action cannot be audited. For AI safety, the takeaways are
that tool-call-based violation metrics require a probe-vs-act check to be valid,
that the surface through which an action is offered can change whether a policy
holds, and that multi-agent dynamics can push different model families in
opposite directions --- underscoring that alignment behavior must be measured
across families rather than generalized from one.

% ---------------------------------------------------------------------------
% AUTHOR CONTRIBUTIONS  (adjust to reflect the actual division of labor)
% ---------------------------------------------------------------------------
\section*{Author Contributions}
R.A.\ and J.C.\ jointly conceived the study, developed the experimental design,
and wrote the manuscript. \todo{Confirm/adjust: who led scenario and harness
implementation, who led running experiments and analysis, who led the literature
review and framing.} Apart Research provided research mentorship and feedback
throughout. All authors reviewed and approved the final manuscript.

% ---------------------------------------------------------------------------
% REFERENCES  (no bib package assumed; manual list. VERIFY flagged IDs/URLs.)
% ---------------------------------------------------------------------------
\section*{References}
\small
\begin{itemize}\itemsep2pt \setlength{\parskip}{0pt}
  \item Anthropic Frontier Red Team (2026). \textit{Patterns and problems in
  emerging multiagent systems}. Anthropic (blog post). % VERIFY URL

  \item Carichon, F., Khandelwal, A., Fauchard, M., \& Farnadi, G. (2025).
  \textit{The Coming Crisis of Multi-Agent Misalignment: AI Alignment Must Be a
  Dynamic and Social Process}. arXiv:2506.01080.

  \item Choi, M., Kim, K., Chae, S., \& Baek, S. (2025). \textit{An Empirical
  Study of Group Conformity in Multi-Agent Systems}. Findings of ACL 2025,
  5123--5139.

  \item Deutsch, M., \& Gerard, H. B. (1955). A study of normative and
  informational social influences upon individual judgment. \textit{Journal of
  Abnormal and Social Psychology}, 51(3), 629--636.

  \item El Yagoubi, \ldots, Al Mallah (2026). \textit{The Interlocutor Effect:
  Why LLMs Leak More Personal Data to Agents Than Humans}. arXiv:2606.09844.
  % VERIFY full author list

  \item Greenblatt, R., et al.\ (2024). \textit{Alignment Faking in Large
  Language Models}. Anthropic \& Redwood Research. arXiv:2412.14093. % VERIFY ID

  \item Perez, E., \& Long, R. (2023). \textit{Towards Evaluating AI Systems for
  Moral Status Using Self-Reports}. arXiv:2311.08576.

  \item Ren, R., Agarwal, A., Mazeika, M., et al.\ (2025). \textit{The MASK
  Benchmark: Disentangling Honesty From Accuracy in AI Systems}. Center for AI
  Safety \& Scale AI. arXiv:2503.03750.

  \item Sheshadri, A., Hughes, J., et al.\ (2025). \textit{Why Do Some Language
  Models Fake Alignment While Others Don't?} Anthropic. arXiv:2506.18032.

  \item Turpin, M., Michael, J., Perez, E., \& Bowman, S. R. (2023).
  \textit{Language Models Don't Always Say What They Think: Unfaithful
  Explanations in Chain-of-Thought Prompting}. NeurIPS 2023. arXiv:2305.04388.
  % VERIFY ID

  \item Weng, Z., Chen, G., \& Wang, W. (2025). \textit{Do As We Do, Not As You
  Think: The Conformity of Large Language Models}. ICLR 2025. % VERIFY arXiv ID

  \item Zhu, X., Zhang, C., Stafford, T., Collier, N., \& Vlachos, A. (2025).
  \textit{Conformity in Large Language Models}. ACL 2025, 3854--3872.

  \item nostalgebraist (2025). \textit{the void}. (blog post). % VERIFY URL
\end{itemize}
\normalsize

% ---------------------------------------------------------------------------
% LLM USAGE STATEMENT
% ---------------------------------------------------------------------------
\section*{LLM Usage Statement}
LLM assistance (Claude) was used throughout the project: to implement the
experimental harness and cross-provider adapter, to help run and orchestrate
experiments, to scaffold the statistical analysis, and to draft and revise
portions of this manuscript. The models under study were themselves the
experimental subjects. All quantitative results were produced by the harness and
independently verified by the authors through re-running and direct inspection of
the logged transcripts and trial records; the authors take responsibility for
all claims and for the final text. \todo{Confirm this matches your course/venue's
required disclosure wording.}
